\documentclass[12pt]{article}
\usepackage[utf8]{inputenc}
\usepackage[T1]{fontenc}
\usepackage{hyperref}
\title{Dangaura Tharu Language Resources for K-12 Curriculum Development}
\author{Compiled by Claude Sonnet 5 (claude-sonnet-5) for Andr\'as Kornai}
\date{}
\begin{document}
\maketitle

\section{Introduction}
Dangaura Tharu (ISO 639-3 thl) is an Indo-Aryan language of Nepal and India. The searches found Bible, audio and sociolinguistic material only. They found no dictionary, grammar or school textbook.

\section{Primary Educational Resources}
None found in the searches.

\section{Linguistic and Reference Materials}
\begin{itemize}
\item Tribhuvan University sociolinguistic study of Dangaura Tharu (2023). Only the contents page was seen; it includes a section on medium of instruction at the primary level: \url{https://portal.tu.edu.np/downloads/Dangaura_Tharu_2023_09_20_20_23_06.pdf}
\end{itemize}

\section{Religious and Community Resources}
\begin{itemize}
\item Scripture Earth (Bible, audio, app): \url{https://www.scriptureearth.org/00eng.php?iso=thl}
\item Scripture Earth index: \url{https://scriptureearth.org/00i-Scripture_Index.php?iso=thl}
\end{itemize}

\section{Digital and Audio Resources}
\begin{itemize}
\item Global Recordings Network, Dangura Tharu: \url{https://globalrecordings.net/en/language/thl}
\item Global Recordings Network, Dasauria Tharu (a possible dialect): \url{https://globalrecordings.net/ja/language/6345}
\end{itemize}

\section{Curriculum Development Context}
Little educational material was found. Any curriculum would have to start from the Bible text and audio, plus original orthography, textbook and dictionary work. The Tribhuvan study may help with decisions on the medium of instruction.
\end{document}